% ======================================================
% 第 61 天
% ======================================================
\setday{61}{方向导数与梯度}

\oneproblem{
    研究 $f(x,y) = \begin{cases} x + y + \dfrac{x^2 + y^2}{\sqrt{x^4 + x^2 + y^2}}, & x^2 + y^2 \neq 0 \\ 0, & x^2 + y^2 = 0 \end{cases}$ 在 $(0,0)$ 处的偏导数和方向导数的存在性.
}

\oneproblem{
    设 $\vec{n}$ 是曲面 $2x^2 + 3y^2 + z^2 = 6$ 在点 $P(1,1,1)$ 处的指向外侧的法向量，求函数 $u = \dfrac{\sqrt{6x^2 + 8y^2}}{z}$ 在点 $P$ 处沿方向 $\vec{n}$ 的方向导数.
}

\oneproblem{
    设 $f(x,y) = \begin{cases} \dfrac{x^2 y^2}{(x^2 + y^2)^{3/2}}, & x^2 + y^2 \neq 0 \\ 0, & x^2 + y^2 = 0 \end{cases}$，试证明函数 $f(x,y)$ 在点 $(0,0)$ 处沿任意方向 $\vec{e}_u = (\cos\alpha, \cos\beta)$ 的方向导数存在，两个偏导数都存在，但不可微.
}

\oneproblem{
    一个蚂蚁在高处觅食的时候，一不小心掉到一块被烤热了的石头上。假设通过建立适当的坐标系后，铁板上温度分布函数为 $f(x,y) = 100 - x^2 - 4y^2$，蚂蚁掉下来的地方位于坐标系中 $(1,-2)$ 的位置，蚂蚁为了免受高温的煎熬，希望沿着温度降低最快的方向逃跑到安全凉爽的地方去！试确定其逃跑的轨迹方程.
}

\oneproblem{
    设函数 $f(x,y,z) = x^2 + y^z$:
    \begin{enumerate}[label=(\arabic*)]
        \item 求函数在点 $M(1,1,1)$ 处沿曲线 $x=t, y=2t^2-1, z=t^3$ 在该点处参变量增大方向的切线方向的方向导数；
        \item 求函数在点 $M(1,1,1)$ 处的梯度与（1）中切线方向的夹角 $\theta$.
    \end{enumerate}
}

\oneproblem{
    设二元函数 $f(x,y)$ 可微，$\vec{l}_1, \vec{l}_2$ 是两个给定的方向，它们之间的夹角为 $\theta \in (0,\pi)$，证明：
    $
    [f_x(x,y)]^2 + [f_y(x,y)]^2 \leq \dfrac{2}{\sin^2\theta}\left[\left(\dfrac{\partial f}{\partial \vec{l}_1}\right)^2 + \left(\dfrac{\partial f}{\partial \vec{l}_2}\right)^2\right].
    $
}

\oneproblem{
    设 $P_0(1,1,-1), P_1(2,-1,0)$ 为空间的两点，求 $u = xyz + e^{xyz}$ 在点 $P_0$ 处沿 $\overrightarrow{P_0P_1}$ 方向的方向导数.
}

\oneproblem{
    设 $\vec{l}_j, j=1,2,\dots,n$ 是平面上点 $P_0$ 处的 $n\geq2$ 各方向向量，相邻两个向量之间的夹角为 $\dfrac{2\pi}{n}$。若函数 $f(x,y)$ 在点 $P_0$ 有连续偏导，证明：$\displaystyle \sum_{j=1}^n \dfrac{\partial f(P_0)}{\partial \vec{l}_i} = 0$.
}

\oneproblem{
    设 $P_0$ 为单位球内一点，从 $P_0$ 向单位球面的切面作垂线，可得垂足 $P$，当切面变动时，$P$ 点的轨迹成一封闭曲面.
    \begin{enumerate}[label=(\arabic*)]
        \item 求此曲面所围成的体积 $V$；
        \item $P_0$ 沿什么方向变动时，此体积的变化率最大？
    \end{enumerate}
}

% ======================================================
% 第 62 天
% ======================================================
\setday{62}{黑塞矩阵及泰勒公式}

\oneproblem{
    设 $f(x,y)$ 具有一阶偏导数，且在任意的 $(x,y)$ 都有 $\dfrac{\partial f(x,y)}{\partial x} > 0, \dfrac{\partial f(x,y)}{\partial y} < 0$，则（ \quad ）
    \begin{multicols}{2}
    \begin{enumerate}[label=(\Alph*)]
        \item $f(0,0) > f(1,1)$
        \item $f(0,0) < f(1,1)$
        \item $f(0,1) > f(1,0)$
        \item $f(0,1) < f(1,0)$
    \end{enumerate}
    \end{multicols}
}

\oneproblem{
    求函数 $f(x,y) = \ln(1+x+y)$ 在点 $(0,0)$ 的三阶带皮亚诺余项的泰勒公式，并求极限
    \[
    \lim_{(x,y)\to(0,0)} \dfrac{\ln(x+y+1)+xy-x-y}{x^2+y^2}.
    \]
}

\oneproblem{
    设函数 $f(x,y)$ 具有一阶连续的偏导数，且 $f(1,0) = f(0,1)$。证明：在单位圆上至少存在两点满足方程 $yf_x - xf_y = 0$.
}

\oneproblem{
    设二元函数 $f(x,y)$ 在点 $(0,0)$ 的某邻域内具有二阶连续偏导数，且 $\displaystyle \lim_{\substack{x\to0 \\ y\to0}} \dfrac{f(x,y)-xy}{(x^2+y^2)^2} = 1$，求 $f(0,0)$ 及 $f(x,y)$ 在点 $(0,0)$ 点的一阶、二阶偏导数的值.
}

\oneproblem{
    设 $f(x,y) = \dfrac{1}{\sqrt{x^2-2xy+1}}$，证明存在 $\theta(0<\theta<1)$，使得
    $
    1-\sqrt{2} = \sqrt{2}(1-3\theta)(1-2\theta+3\theta^2)^{-3/2}.
    $
}

\oneproblem{
    设 $n$ 元函数 $f(\mathbf{x})$ 在点 $\mathbf{x}_0$ 处具有二阶连续偏导数，且 $\nabla f(\mathbf{x}_0) = 0$，记 $\mathbf{H}(\mathbf{x}_0)$ 为 $f(\mathbf{x})$ 在点 $\mathbf{x}_0$ 处的黑塞矩阵。证明：
    \begin{enumerate}[label=(\arabic*)]
        \item 如果 $\mathbf{H}(\mathbf{x}_0)$ 正定，则存在 $U_0(\mathbf{x}_0)$，使得 $\forall \mathbf{x} \in U_0(\mathbf{x}_0)$，有 $f(\mathbf{x}_0) < f(\mathbf{x})$；
        \item 如果 $\mathbf{H}(\mathbf{x}_0)$ 负定，则存在 $U_0(\mathbf{x}_0)$，使得 $\forall \mathbf{x} \in U_0(\mathbf{x}_0)$，有 $f(\mathbf{x}_0) > f(\mathbf{x})$.
    \end{enumerate}
    \vspace{0.5em}
    \noindent\footnotesize{[说明]: 矩阵 $\mathbf{A}$ 正定是指任给非零向量 $\mathbf{x}$，有 $\mathbf{x}^T\mathbf{A}\mathbf{x} > 0$；负定则是 $\mathbf{x}^T\mathbf{A}\mathbf{x} < 0$.}
}

% ======================================================
% 第 63 天
% ======================================================
\setday{63}{多元函数的极值与条件极值}

\oneproblem{
    设可微函数 $f(x,y)$ 在点 $(x_0,y_0)$ 取得极小值，则下列结论正确的是（ \quad ）
    \begin{multicols}{2}
    \begin{enumerate}[label=(\Alph*)]
        \item $f(x_0,y)$ 在 $y=y_0$ 处的导数等于零
        \item $f(x_0,y)$ 在 $y=y_0$ 处的导数大于零
        \item $f(x_0,y)$ 在 $y=y_0$ 处的导数小于零
        \item $f(x_0,y)$ 在 $y=y_0$ 处的导数不存在
    \end{enumerate}
    \end{multicols}
}

\oneproblem{
    已知函数 $f(x,y)$ 在点 $(0,0)$ 的某个邻域内连续，且 $\displaystyle \lim_{\substack{x\to0 \\ y\to0}} \dfrac{f(x,y)-xy}{(x^2+y^2)^2} = 1$，则（ \quad ）
    \begin{enumerate}[label=(\Alph*)]
        \item 点 $(0,0)$ 不是 $f(x,y)$ 的极值点
        \item 点 $(0,0)$ 是 $f(x,y)$ 的极大值点
        \item 点 $(0,0)$ 是 $f(x,y)$ 的极小值点
        \item 根据所给条件无法判断点 $(0,0)$ 是否为 $f(x,y)$ 的极值点
    \end{enumerate}
}

\oneproblem{
    设函数 $z = f(x,y)$ 的全微分为 $\mathrm{d}z = x\mathrm{d}x + y\mathrm{d}y$，则点 $(0,0)$（ \quad ）
    \begin{multicols}{2}
    \begin{enumerate}[label=(\Alph*)]
        \item 不是 $f(x,y)$ 的连续点
        \item 不是 $f(x,y)$ 的极值点
        \item 是 $f(x,y)$ 的极大值点
        \item 是 $f(x,y)$ 的极小值点
    \end{enumerate}
    \end{multicols}
}

\oneproblem{
    设函数 $u(x,y)$ 在有界闭区域 $D$ 上连续，在 $D$ 的内部具有二阶连续偏导数，且满足 $\dfrac{\partial^2 u}{\partial x \partial y} \neq 0$ 及 $\dfrac{\partial^2 u}{\partial x^2} + \dfrac{\partial^2 u}{\partial y^2} = 0$，则（ \quad ）
    \begin{enumerate}[label=(\Alph*)]
        \item $u(x,y)$ 的最大值和最小值都在 $D$ 的边界上取得
        \item $u(x,y)$ 的最大值和最小值都在 $D$ 的内部取得
        \item $u(x,y)$ 的最大值在 $D$ 的内部取得，最小值在 $D$ 的边界上取得
        \item $u(x,y)$ 的最小值在 $D$ 的内部取得，最大值在 $D$ 的边界上取得
    \end{enumerate}
}

\oneproblem{
    设 $f(x,y)$ 与 $\varphi(x,y)$ 均为可微函数，且 $\varphi'_y(x,y) \neq 0$，已知 $(x_0,y_0)$ 是 $f(x,y)$ 在约束条件 $\varphi(x,y)=0$ 下的一个极值点，下列选项正确的是（ \quad ）
    \begin{multicols}{2}
    \begin{enumerate}[label=(\Alph*)]
        \item 若 $f'_x(x_0,y_0) = 0$，则 $f'_y(x_0,y_0) = 0$
        \item 若 $f'_x(x_0,y_0) = 0$，则 $f'_y(x_0,y_0) \neq 0$
        \item 若 $f'_x(x_0,y_0) \neq 0$，则 $f'_y(x_0,y_0) = 0$
        \item 若 $f'_x(x_0,y_0) \neq 0$，则 $f'_y(x_0,y_0) \neq 0$
    \end{enumerate}
    \end{multicols}
}

\oneproblem{
    设函数 $f(x),g(x)$ 均有二阶连续导数，满足 $f(0) > 0, g(0) < 0$ 且 $f'(0) = g'(0) = 0$，则函数 $z = f(x)g(y)$ 在点 $(0,0)$ 处取得极小值的一个充分条件是（ \quad ）
    \begin{multicols}{2}
    \begin{enumerate}[label=(\Alph*)]
        \item $f''(0) < 0, g''(0) > 0$
        \item $f''(0) < 0, g''(0) < 0$
        \item $f''(0) > 0, g''(0) > 0$
        \item $f''(0) > 0, g''(0) < 0$
    \end{enumerate}
    \end{multicols}
}

\oneproblem{
    当 $x\geq0, y\geq0$ 时，$x^2 + y^2 \leq k e^{x+y}$ 恒成立，试求 $k$ 的取值范围.
}

\oneproblem{
    求二元函数 $f(x,y) = x^2 y (4 - x - y)$ 在由直线 $x+y=6$，$x$ 轴和 $y$ 轴所围成的闭区域 $D$ 上的极值、最大值与最小值.
}

\oneproblem{
    设有一小山，取它的底面所在的平面为 $xOy$ 坐标面，其底部所占的区域为 $D = \{(x,y) \mid x^2 + y^2 - xy \leq 75\}$，小山的高度函数为 $h(x,y) = 75 - x^2 - y^2 + xy$.
    \begin{enumerate}[label=(\arabic*)]
        \item 设 $M(x_0,y_0)$ 为区域 $D$ 上一点，问 $h(x,y)$ 在该点沿平面上什么方向的方向导数最大？若记此方向导数的最大值为 $g(x_0,y_0)$，试写出 $g(x_0,y_0)$ 的表达式.
        \item 现欲利用此小山开展攀岩活动，为此需要在山脚寻找一上山坡度最大的点作为攀登的起点，也就是说，要在 $D$ 的边界线 $x^2 + y^2 - xy = 75$ 上找出使 (1) 中的 $g(x,y)$ 达到最大值的点。试确定攀登起点的位置.
    \end{enumerate}
}

\oneproblem{
    已知函数 $f(x,y) = x + y + xy$，曲线 $C: x^2 + y^2 + xy = 3$，求 $f(x,y)$ 在曲线 $C$ 上的最大方向导数.
}

\oneproblem{
    已知函数 $z = z(x,y)$ 由方程 $(x^2 + y^2)z + \ln z + 2(x + y + 1) = 0$ 确定，求 $z = z(x,y)$ 的极值.
}

\oneproblem{
    已知函数 $f(x,y)$ 满足 $f''_{xy}(x,y) = 2(y+1)e^x$，$f'_x(x,0) = (x+1)e^x$，$f(0,y) = y^2 + 2y$，求 $f(x,y)$ 的极值.
}

\oneproblem{
    将长为 $2m$ 的铁丝分成三段，分别围成圆、正方形与正三角形，三个图形的面积之和是否存在最小值，如果存在，求出最小值.
}

\oneproblem{
    已知可微函数 $f(u,v)$ 满足
    $
    \dfrac{\partial f(u,v)}{\partial u} - \dfrac{\partial f(u,v)}{\partial v} = 2(u-v)e^{-(u+v)}
    $
    且 $f(u,0) = u^2 e^{-u}$.
    \begin{enumerate}[label=(\arabic*)]
        \item 记 $g(x,y) = f(x, y-x)$，求 $\dfrac{\partial g(x,y)}{\partial x}$.
        \item 求 $f(u,v)$ 的表达式和极值.
    \end{enumerate}
}

\oneproblem{
    对于 $\triangle ABC$，求 $3\sin A + 4\sin B + 18\sin C$ 的最大值.
}

\oneproblem{
    设二元函数 $f(x,y)$ 在平面上有连续的二阶偏导数，对任意角度 $\alpha$，定义一元函数 $g_\alpha(t) = f(t\cos\alpha, t\sin\alpha)$，若对任何 $\alpha$ 都有 $\dfrac{\mathrm{d}g_\alpha(0)}{\mathrm{d}t} = 0$ 且 $\dfrac{\mathrm{d}^2 g_\alpha(0)}{\mathrm{d}t^2} > 0$，证明：$f(0,0)$ 是 $f(x,y)$ 的极小值.
}

\oneproblem{
    试求函数 $u = x_1 + \dfrac{x_2}{x_1} + \dfrac{x_3}{x_2} + \dfrac{2}{x_3} \ (x_i > 0, i=1,2,3)$ 的所有极值点.
}

\oneproblem{
    设 $\Sigma_1: \dfrac{x^2}{a^2} + \dfrac{y^2}{b^2} + \dfrac{z^2}{c^2} = 1$，其中 $a>b>c>0$，$\Sigma_2: z^2 = x^2 + y^2$，$\Gamma$ 为 $\Sigma_1$ 和 $\Sigma_2$ 的交线。求椭球面 $\Sigma_1$ 在 $\Gamma$ 上各点的切平面到原点距离的最大值和最小值.
}

\oneproblem{
    设 $D = \{(x,y): x^2 + y^2 < 1\}$，$f(x,y)$ 在 $D$ 内连续，$g(x,y)$ 在 $D$ 内连续有界，且满足条件：
    \begin{enumerate}[label=(\arabic*)]
        \item 当 $x^2 + y^2 \to 1$ 时，$f(x,y) \to +\infty$；
        \item 在 $D$ 内 $f$ 与 $g$ 有二阶偏导数，$\dfrac{\partial^2 f}{\partial x^2} + \dfrac{\partial^2 f}{\partial y^2} = e^f$ 和 $\dfrac{\partial^2 g}{\partial x^2} + \dfrac{\partial^2 g}{\partial y^2} \geq e^g$.
    \end{enumerate}
    证明：$f(x,y) \geq g(x,y)$ 在 $D$ 内处处成立.
}

% ======================================================
% 第 64 天
% ======================================================
\setday{64}{二重积分的概念和性质}

\oneproblem{
    设 $D$ 是 $xOy$ 平面上以 $(1,1),(-1,1)$ 和 $(-1,-1)$ 为顶点的三角区域，$D_1$ 是 $D$ 在第一象限的部分，则 $\displaystyle \iint_D (xy + \cos x \sin y)\,\dd x \dd y$ 等于（ \quad ）
    \begin{multicols}{2}
    \begin{enumerate}[label=(\Alph*)]
        \item $2\iint_{D_1} \cos x \sin y\,\dd x \dd y$
        \item $2\iint_{D_1} xy\,\dd x \dd y$
        \item $4\iint_{D_1} (xy + \cos x \sin y)\,\dd x \dd y$
        \item $0$
    \end{enumerate}
    \end{multicols}
}

\oneproblem{
    设 $I_1 = \iint_D \cos\sqrt{x^2+y^2}\,d\sigma,\ I_2 = \iint_D \cos(x^2+y^2)\,d\sigma,\ I_3 = \iint_D \cos(x^2+y^2)^2\,d\sigma$，其中 $D = \{(x,y)|x^2+y^2\le1\}$，则（ \quad ）
    \begin{multicols}{2}
    \begin{enumerate}[label=(\Alph*)]
        \item $I_3>I_2>I_1$
        \item $I_1>I_2>I_3$
        \item $I_2>I_1>I_3$
        \item $I_3>I_1>I_2$
    \end{enumerate}
    \end{multicols}
}

\oneproblem{
    设区域 $D = \{(x,y)|x^2+y^2\le4,x\ge0,y\ge0\}$，$f(x)$ 为 $D$ 上的正值连续函数，$a,b$ 为常数，则 $\displaystyle \iint_D \dfrac{a\sqrt{f(x)}+b\sqrt{f(y)}}{\sqrt{f(x)}+\sqrt{f(y)}}\,d\sigma =$（ \quad ）
    \begin{multicols}{2}
    \begin{enumerate}[label=(\Alph*)]
        \item $ab\pi$
        \item $\dfrac{ab}{2}\pi$
        \item $(a+b)\pi$
        \item $\dfrac{a+b}{2}\pi$
    \end{enumerate}
    \end{multicols}
}

\oneproblem{
    设 $f(x)$ 是连续的奇函数，$g(x)$ 是连续的偶函数，区域 $D = \{(x,y)|0\le x\le1,-\sqrt{x}\le y\le\sqrt{x}\}$，则下列结论正确的是（ \quad ）
    \begin{multicols}{2}
    \begin{enumerate}[label=(\Alph*)]
        \item $\iint_D f(y)g(x)\,\dd x \dd y=0$
        \item $\iint_D f(x)g(y)\,\dd x \dd y=0$
        \item $\iint_D [f(x)+g(y)]\,\dd x \dd y=0$
        \item $\iint_D [f(y)+g(x)]\,\dd x \dd y=0$
    \end{enumerate}
    \end{multicols}
}

\oneproblem{
    如图所示，正方形 $\{(x,y)|\ |x|\le1,\ |y|\le1\}$ 被其对角线划分为四个区域 $D_k\ (k=1,2,3,4)$，$I_k = \iint_{D_k} y\cos x\,\dd x \dd y$，则 $\max_{1\le k\le4}\{I_k\} =$（ \quad ）
    \begin{multicols}{2}
    \begin{enumerate}[label=(\Alph*)]
        \item $I_1$
        \item $I_2$
        \item $I_3$
        \item $I_4$
    \end{enumerate}
    \end{multicols}
}

\oneproblem{
    $\displaystyle \lim_{n\to\infty} \sum_{i=1}^n\sum_{j=1}^n \dfrac{n}{(n+i)(n^2+j^2)} =$（ \quad ）
    \begin{multicols}{2}
    \begin{enumerate}[label=(\Alph*)]
        \item $\int_0^1 dx\int_0^x \dfrac{1}{(1+x)(1+y^2)}\,dy$
        \item $\int_0^1 dx\int_0^x \dfrac{1}{(1+x)(1+y)}\,dy$
        \item $\int_0^1 dx\int_0^1 \dfrac{1}{(1+x)(1+y)}\,dy$
        \item $\int_0^1 dx\int_0^1 \dfrac{1}{(1+x)(1+y^2)}\,dy$
    \end{enumerate}
    \end{multicols}
}

\oneproblem{
    设 $D_k$ 是圆域 $D=\{(x,y)|\ x^2+y^2\le1\}$ 的第 $k$ 象限的部分，记 $I_k=\iint_{D_k}(y-x)\,\dd x \dd y\ (k=1,2,3,4)$，则（ \quad ）
    \begin{multicols}{2}
    \begin{enumerate}[label=(\Alph*)]
        \item $I_1>0$
        \item $I_2>0$
        \item $I_3>0$
        \item $I_4>0$
    \end{enumerate}
    \end{multicols}
}

\oneproblem{
    设 $J_i = \iint_{D_i} \sqrt[3]{x-y}\,\dd x \dd y\ (i=1,2,3)$，其中$D_1=\{(x,y)|\ 0\le x\le1,0\le y\le1\},$\\$\quad D_2=\{(x,y)|\ 0\le x\le1,0\le y\le\sqrt{x}\},\quad D_3=\{(x,y)|\ 0\le x\le1,x^2\le y\le1\},
    $
    则（ \quad ）
    \begin{multicols}{2}
    \begin{enumerate}[label=(\Alph*)]
        \item $J_1<J_2<J_3$
        \item $J_3<J_1<J_2$
        \item $J_2<J_3<J_1$
        \item $J_2<J_1<J_3$
    \end{enumerate}
    \end{multicols}
}

\oneproblem{
    已知积分区域 $D=\{(x,y)|\ |x|+|y|\le\dfrac{\pi}{2}\}$，记
    $
    I_1=\iint_D \sqrt{x^2+y^2}\,\dd x \dd y,\quad I_2=\iint_D \sin\sqrt{x^2+y^2}\,\dd x \dd y,\quad I_3=\iint_D (1-\cos\sqrt{x^2+y^2})\,\dd x \dd y,
    $
    比较 $I_1,I_2,I_3$ 的大小（ \quad ）
    \begin{multicols}{2}
    \begin{enumerate}[label=(\Alph*)]
        \item $I_3<I_2<I_1$
        \item $I_1<I_2<I_3$
        \item $I_2<I_1<I_3$
        \item $I_2<I_3<I_1$
    \end{enumerate}
    \end{multicols}
}

\oneproblem{
    设平面区域 $D=\{(x,y)|\ x^3\le y\le1,-1\le x\le1\}$，$f(x)$ 是定义在 $[-a,a]\ (a\ge1)$ 上的任意连续函数，计算二重积分
    $
    I = \iint_D 2y[(x+1)f(x)+(x-1)f(-x)]\,\dd x \dd y.
    $
}

\oneproblem{
    判断积分 $I = \iint_D \sqrt[3]{1-x^2-y^2}\,\dd x \dd y$ 的正负号，其中 $D:\ x^2+y^2\le4$.
}

\oneproblem{
    计算 $\displaystyle \lim_{n\to\infty} \sum_{i=1}^n\sum_{j=1}^{2n} \dfrac{2}{n^2}\left[\dfrac{2i+j}{n}\right]$，这里 $[x]$ 是不超过 $x$ 的最大整数.
}

\oneproblem{
    设 $D$ 是平面上由光滑闭曲线围成的有界区域，其面积为 $A>0$，函数 $f(x,y)$ 在该区域及其边界上连续，函数 $f(x,y)$ 在 $D$ 上连续且 $f(x,y)>0$。记 $J_n = \left(\dfrac{1}{A}\iint_D f^{1/n}(x,y)\,d\sigma\right)^n$，计算 $\displaystyle \lim_{n\to+\infty} J_n$.
    \vspace{0.5em}
    \noindent\footnotesize{[提示] 假设满足题设条件的积分有 $G'(t)=\dfrac{d}{dt}\iint_D g(x,y,t)\,d\sigma=\iint_D \dfrac{dg(x,y,t)}{dt}\,d\sigma$.}
}

\oneproblem{
    设 $D=\{(x,y)|\ 0\le x\le1,0\le y\le1\}$，$I=\iint_D f(x,y)\,\dd x \dd y$，其中 $f(x,y)$ 在 $D$ 上有连续二阶偏导数，若对任何 $x,y$ 有 $f(0,y)=f(x,0)=0$ 且 $\dfrac{\partial^2 f}{\partial x\partial y}\le A$。证明：$I\le\dfrac{A}{4}$.
}

\oneproblem{
    设 $f:[-1,1]\to R$ 为偶函数，$f$ 在 $[0,1]$ 上单调递增，又设 $g$ 是 $[-1,1]$ 上的凸函数，即对任意 $x,y\in[-1,1]$ 及 $t\in(0,1)$ 有 $g(tx+(1-t)y)\le tg(x)+(1-t)g(y)$。求证：$2\int_{-1}^1 f(x)g(x)\,dx\ge\int_{-1}^1 f(x)\,dx\int_{-1}^1 g(x)\,dx$.
}

\oneproblem{
    设函数 $f(x)$ 在 $[0,1]$ 上连续，满足对任意 $x\in[0,1]$，$\int_{x^2}^x f(t)\,dt\ge\dfrac{x^2-x^4}{2}$。证明：$\int_0^1 f^2(x)\,dx\ge\dfrac{1}{10}$.
}

\oneproblem{
    设 $f(x,y)$ 是定义在区域 $0\le x\le1,0\le y\le1$ 上的二元函数，$f(0,0)=0$，且在点 $(0,0)$ 处 $f(x,y)$ 可微分。证明：
    $
    \lim_{x\to0^+}\dfrac{\int_0^{x^2} dt\int_x^{\sqrt{t}} f(t,u)\,du}{1-e^{-\frac{x^4}{4}}} = -\left.\dfrac{\partial f}{\partial y}\right|_{(0,0)}.
    $
}

% ======================================================
% 第 65 天
% ======================================================
\setday{65}{直角坐标系下二重积分的计算}

\oneproblem{
    计算下列极限
    \begin{enumerate}[label=(\arabic*)]
        \item $\displaystyle \lim_{n\to\infty} \dfrac{1}{n}\sum_{i=1}^n\sum_{j=1}^n \dfrac{i+j}{i^2+j^2}$
        \item $\displaystyle \lim_{n\to\infty} \dfrac{\pi}{2n^4}\sum_{i=1}^n i^2\sin\dfrac{i\pi}{2n}$
    \end{enumerate}
}

\oneproblem{
    已知函数 $f(t)=\int_1^{t^2} dx\int_{\sqrt{x}}^t \sin\dfrac{x}{y}\,dy$，求 $f'\left(\dfrac{\pi}{2}\right)$.
}

\oneproblem{
    计算下列累次积分.
    \begin{enumerate}[label=(\arabic*)]
        \item $I=\int_0^{2\pi} x\,dx\int_x^{2\pi} \dfrac{\sin^2 t}{t^2}\,dt$
        \item $I=\int_1^2 dx\int_{\sqrt{x}}^x \sin\dfrac{\pi x}{2y}\,dy + \int_2^4 dx\int_{\sqrt{x}}^2 \sin\dfrac{\pi x}{2y}\,dy$
    \end{enumerate}
}

\oneproblem{
    计算积分 $\displaystyle \iint_D \dfrac{y^3}{(1+x^2+y^4)^2}\,\dd x \dd y$，其中 $D$ 是第一象限中曲线 $y=\sqrt{x}$ 与 $x$ 轴边界的无界区域。
}

\oneproblem{
    设平面区域 $D$ 由曲线 $\begin{cases} x=t-\sin t \\ y=1-\cos t \end{cases}\ (0\le t\le2\pi)$ 与 $x$ 轴围成，计算二重积分 $\displaystyle \iint_D (x+2y)\,\dd x \dd y$.
}

\oneproblem{
    计算二重积分 $\displaystyle \iint_D (x+y)\,\dd x \dd y$，其中 $D=\{(x,y)|\ x^2+y^2\le x+y+1\}$.
}

\oneproblem{
    设 $f(x,y)=\begin{cases} x^2y, & 1\le x\le2,0\le y\le x \\ 0, & \text{其他} \end{cases}$，求 $\displaystyle \iint_D f(x,y)\,\dd x \dd y$，其中 $D=\{(x,y)|\ x^2+y^2\ge2x\}$.
}

\oneproblem{
    设 $f(x,y)$ 连续，且 $f(x,y)=xy+\iint_D f(u,v)\,dudv$，其中 $D$ 是由 $y=0,\ y=x^2,\ x=1$ 所围成的区域，试求 $f(x,y)$ 的表达式.
}

\oneproblem{
    设函数 $f(x)$ 在 $[0,1]$ 有连续导数，$f(0)=1$，且满足 $\displaystyle \iint_{D_t} f'(x+y)\,\dd x \dd y=\iint_{D_t} f(t)\,\dd x \dd y$，其中 $D_t=\{(x,y)|\ 0\le y\le t-x,0\le x\le t\}\ (0<t\le1)$，求 $f(x)$ 的表达式.
}

\oneproblem{
    已知 $f(x)$ 在 $[0,\dfrac{3\pi}{2}]$ 上连续，在 $(0,\dfrac{3\pi}{2})$ 内是函数 $\dfrac{\cos x}{2x-3\pi}$ 的一个原函数，且 $f(0)=0$.
    \begin{enumerate}[label=(\arabic*)]
        \item 求 $f(x)$ 在区间 $[0,\dfrac{3\pi}{2}]$ 上的平均值
        \item 证明 $f(x)$ 在区间 $(0,\dfrac{3\pi}{2})$ 内存在唯一零点
    \end{enumerate}
}

\oneproblem{
    已知函数 $f(x,y)$ 具有二阶连续偏导数，且 $f(1,y)=0,\ f(x,1)=0,\ \displaystyle \iint_D f(x,y)\,\dd x \dd y=a$，其中 $D=\{(x,y)|\ 0\le x\le1,0\le y\le1\}$，计算二重积分 $I=\iint_D xyf''_{xy}(x,y)\,\dd x \dd y$.
}

% ======================================================
% 第 66 天
% ======================================================
\setday{66}{极坐标系下二重积分的计算}

\oneproblem{
    累次积分 $\displaystyle \int_0^{\frac{\pi}{2}} d\theta\int_0^{\cos\theta} f(r\cos\theta,r\sin\theta)r\,dr$（ \quad ）
    \begin{multicols}{2}
    \begin{enumerate}[label=(\Alph*)]
        \item $\int_0^1 dy\int_0^{\sqrt{y-y^2}} f(x,y)\,dx$
        \item $\int_0^1 dy\int_0^{\sqrt{1-y^2}} f(x,y)\,dx$
        \item $\int_0^1 dx\int_0^1 f(x,y)\,dy$
        \item $\int_0^1 dx\int_0^{\sqrt{x-x^2}} f(x,y)\,dy$
    \end{enumerate}
    \end{multicols}
}

\oneproblem{
    设函数 $f$ 连续，若 $F(u,v)=\iint_{D_{uv}} \dfrac{f(x^2+y^2)}{\sqrt{x^2+y^2}}\,\dd x \dd y$，其中 $D_{uv}:\ x^2+y^2=1,\ x^2+y^2=u^2,\ y=0,\ y=x\tan v\ (u>1,v>0)$，则 $\dfrac{\partial F}{\partial u}=$（ \quad ）
    \begin{multicols}{2}
    \begin{enumerate}[label=(\Alph*)]
        \item $vf(u^2)$
        \item $\dfrac{v}{u}f(u^2)$
        \item $vf(u)$
        \item $\dfrac{v}{u}f(u)$
    \end{enumerate}
    \end{multicols}
}

\oneproblem{
    计算下列二重积分.
    \begin{enumerate}[label=(\arabic*),itemsep=8em]
        \item $\displaystyle \iint_D (\sin x^2\cos y^2 + x\sqrt{x^2+y^2})\,\dd x \dd y$，其中 $D:\ x^2+y^2\le\pi$
        \item $\displaystyle \iint_D (x+y^2)e^{-(x^2+y^2-4)}\,\dd x \dd y$，其中 $D:\ 1\le x^2+y^2\le4$
        \item $\displaystyle \iint_D \dfrac{x\sin(\pi\sqrt{x^2+y^2})}{x+y}\,\dd x \dd y$，其中 $D:\ 1\le x^2+y^2\le4,x\ge0,y\ge0$
        \item $\displaystyle \iint_D \dfrac{x+y}{\sqrt{x^2+y^2}}\,\dd x \dd y$，其中 $D:\ |x|\le y,(x^2+y^2)^3\le y^4$
        \item $\displaystyle \iint_D r^2\sin\theta\sqrt{1-r^2\cos2\theta}\,dr d\theta$，其中 $D:\ 0\le r\le\sec\theta,0\le\theta\le\dfrac{\pi}{4}$
    \end{enumerate}
}

\oneproblem{
    设 $D:\ x^2+y^2\le1$，证明：$\dfrac{61}{165}\pi\le\iint_D \sin\sqrt{(x^2+y^2)^3}\,d\sigma\le\dfrac{2}{5}\pi$.
}

\oneproblem{
    设函数 $f(x)$ 在 $[0,+\infty)$ 上连续，且满足方程 $f(t)=e^{4\pi t^2}+\displaystyle \iint_{x^2+y^2\le4t^2} f\left(\dfrac{1}{2}\sqrt{x^2+y^2}\right)\,\dd x \dd y$，求 $f(t)$.
}

\oneproblem{
    设函数 $z=f(x,y)$ 满足方程 $F(u,v)=0$，其中 $u=x+az,\ v=y+bz,\ a,b$ 为常数，$F$ 可微，且 $aF_u+bF_v\neq0$，求积分 $I=\displaystyle \iint_{x^2+y^2\le1} e^{-(x^2+y^2)}(az_x+bz_y)\,\dd x \dd y$.
}

\oneproblem{
    设区域 $D=\{(x,y)|\ x^2+y^2\le1,y\ge0\}$，连续函数 $f(x,y)$ 满足 $f(x,y)=y\sqrt{1-x^2}+x\displaystyle \iint_D f(x,y)\,\dd x \dd y$，计算 $\displaystyle \iint_D xf(x,y)\,\dd x \dd y$.
}

\oneproblem{
    设 $f(x,y)$ 在 $x^2+y^2\le1$ 上有连续的二阶导数，$f_{xx}^2+2f_{xy}^2+f_{yy}^2\le M$。若 $f(0,0)=f_x(0,0)=f_y(0,0)=0$，证明：$\left|\displaystyle \iint_{x^2+y^2\le1} f(x,y)\,\dd x \dd y\right|\le\dfrac{\pi\sqrt{M}}{4}$.
}

\oneproblem{
    设函数 $f(x,y)$ 在单位圆域上有连续的偏导数，且在边界上的值恒为零。证明：
    $
    \lim_{\varepsilon\to0^+} \dfrac{-1}{2\pi}\iint_D \dfrac{xf_x'+yf_y'}{x^2+y^2}\,\dd x \dd y = f(0,0),
    $
    其中 $D$ 为圆环域 $\varepsilon^2\le x^2+y^2\le1$.
}

\oneproblem{
    设 $f(x,y)\ge0$ 在 $D:\ x^2+y^2\le a^2$ 上有连续的一阶偏导数，边界上取值为零。证明不等式：
    $
    \left|\displaystyle \iint_D f(x,y)\,\dd x \dd y\right|\le\dfrac{\pi a^3}{3}\max_{(x,y)\in D}\sqrt{\left(\dfrac{\partial f}{\partial x}\right)^2+\left(\dfrac{\partial f}{\partial y}\right)^2}.
    $
}

% ======================================================
% 第 67 天
% ======================================================
\setday{67}{三重积分的概念 直角坐标系下三重积分的计算}

\oneproblem{
    设有空间区域
    $
    \Omega_1: x^2+y^2+z^2\le R^2, z\ge0\quad\text{及}\quad\Omega_2: x^2+y^2+z^2\le R^2, x\ge0, y\ge0, z\ge0,
    $
    则（ \quad ）
    \begin{multicols}{2}
    \begin{enumerate}[label=(\Alph*)]
        \item $\displaystyle \iiint_{\Omega_1} x\,dV = 4\iiint_{\Omega_2} x\,dV$
        \item $\displaystyle \iiint_{\Omega_1} y\,dV = 4\iiint_{\Omega_2} y\,dV$
        \item $\displaystyle \iiint_{\Omega_1} z\,dV = 4\iiint_{\Omega_2} z\,dV$
        \item $\displaystyle \iiint_{\Omega_1} xyz\,dV = 4\iiint_{\Omega_2} xyz\,dV$
    \end{enumerate}
    \end{multicols}
}

\oneproblem{
    设 $\Omega$ 是由平面 $x+y+z=1$ 与三个坐标平面所围成的空间区域，计算三重积分 $\displaystyle \iiint_{\Omega} (x+2y+3z)\,dV$.
}

\oneproblem{
    记曲面 $z^2=x^2+y^2$ 和 $z=\sqrt{4-x^2-y^2}$ 围成的空间区域为 $V$，计算三重积分 $\displaystyle \iiint_V z\,\dd x \dd ydz$.
}

\oneproblem{
    设 $\Omega$ 为 $x^2+y^2+z^2\le1$。证明：
    $
    \frac{4\sqrt[3]{2}}{3}\pi \le \iiint_{\Omega} \sqrt[3]{x+2y-2z+5}\,dV \le \frac{8}{3}\pi.
    $
}

\oneproblem{
    求 $I=\displaystyle \iiint_{\Omega} (x^2+y^2+z)\,dV$，其中 $\Omega$ 是由曲线 $\begin{cases} y^2=2z \\ x=0 \end{cases}$ 绕 $z$ 轴旋转一周而成的曲面与平面 $z=4$ 所围成的立体.
}

\oneproblem{
    计算三重积分 $I=\displaystyle \iiint_{\Omega} (x+y+z)^2\,dV$，其中 $\Omega: x^2+y^2+z^2\le1$.
}

\oneproblem{
    计算 $I=\displaystyle \iiint_{\Omega} (x+y)\,dV$，其中 $\Omega$ 由曲面 $y=-\sqrt{1-x^2-z^2},\ x^2+z^2=1,\ y=1$ 所围成.
}

\oneproblem{
    计算三次积分 $I=\displaystyle \int_0^1 dx\int_0^{1-x} dy\int_0^{1-x-y} (1-y)e^{-(1-y-z)^2}\,dz$.
}

\oneproblem{
    设 $f(x)$ 为 $[0,1]$ 上的连续函数，证明：
    $
    \int_0^1 dx\int_x^1 dy\int_x^y f(x)f(y)f(z)\,dz = \frac{1}{6}\left[\int_0^1 f(t)\,dt\right]^3.
    $
}

\oneproblem{
    计算三重积分 $\displaystyle \iiint_{\Omega} \frac{\dd x \dd ydz}{(1+x^2+y^2+z^2)^2}$，其中 $\Omega: 0\le x\le1, 0\le y\le1, 0\le z\le1$.
}

% ======================================================
% 第 68 天
% ======================================================
\setday{68}{柱面坐标系与球面坐标系下三重积分的计算}

\oneproblem{
    设 $\Omega$ 为由曲面 $z=\sqrt{x^2+y^2},\ z=\sqrt{1-x^2-y^2}$ 所围成的立体，试分别用柱面坐标与球面坐标方法计算三重积分 $I=\displaystyle \iiint_{\Omega} \sqrt{x^2+y^2}\,dV$.
}

\oneproblem{
    计算三重积分 $I=\displaystyle \iiint_{\Omega} (x+z)\,dV$，其中 $\Omega$ 是由曲面 $z=\sqrt{x^2+y^2}$ 与 $z=\sqrt{1-x^2-y^2}$ 所围成的区域.
}

\oneproblem{
    求曲面 $(x^2+y^2)^2+z^4=y$ 所围立体的体积.
}

\oneproblem{
    将下列直角坐标系下的三次积分
    $
    \int_0^1 dy\int_{-\sqrt{y-y^2}}^{\sqrt{y-y^2}} dx\int_0^{\sqrt[3]{x^2+y^2}} f\left(\sqrt{x^2+y^2+z^2}\right)\,dz
    $
    分别转化为柱面坐标系和球面坐标系下的三次积分.
}

\oneproblem{
    设 $f$ 在 $[0,1]$ 上连续，$f(0)=0,\ f'(0)=1$。求极限
    $
    \lim_{t\to0^+} \frac{1}{t^4} \iiint_{x^2+y^2+z^2\le t^2} f\left(\sqrt{x^2+y^2+z^2}\right)\,\dd x \dd ydz.
    $
}

\oneproblem{
    设函数 $f(x)$ 连续且恒大于零，
    $
    F(t) = \frac{\iiint_{\Omega(t)} f(x^2+y^2+z^2)\,dv}{\iint_{D(t)} f(x^2+y^2)\,d\sigma},\quad G(t) = \frac{\iint_{D(t)} f(x^2+y^2)\,d\sigma}{\int_{-t}^t f(x^2)\,dx},
    $
    其中 $\Omega(t)=\{(x,y,z)|\ x^2+y^2+z^2\le t^2\},\ D(t)=\{(x,y)|\ x^2+y^2\le t^2\}$.
    \begin{enumerate}[label=(\arabic*)]
        \item 讨论 $F(t)$ 在区间 $(0,+\infty)$ 内的单调性.
        \item 证明当 $t>0$ 时，$F(t)>\dfrac{2}{\pi}G(t)$.
    \end{enumerate}
}

\oneproblem{
    设 $f(x)$ 为连续函数，$\Omega$ 是由抛物面 $z=x^2+y^2$ 和球面 $x^2+y^2+z^2=t^2\ (t>0)$ 所围成起来的部分。定义 $F(t)=\displaystyle \iiint_{\Omega} f(x^2+y^2+z^2)\,dV$，求 $F'(t)$.
}

\oneproblem{
    设 $\Omega: x^2+y^2+z^2\le1$，证明积分等式
    $
    \iiint_{\Omega} f(z)\,dz = \pi\int_{-1}^1 f(u)(1-u^2)\,du.
    $
}

\oneproblem{
    计算三重积分 $\displaystyle \iiint_{\Omega} \frac{xyz}{x^2+y^2}\,\dd x \dd ydz$，其中 $\Omega$ 是由曲面 $(x^2+y^2+z^2)^2=2xy$ 围成的区域在第一卦限部分.
}

\oneproblem{
    计算三重积分 $\displaystyle \iiint_{(V)} (x^2+y^2)\,dV$，其中 $(V)$ 是由 $x^2+y^2+(z-2)^2\ge4,\ x^2+y^2+(z-1)^2\le9$ 及 $z\ge0$ 所围成的空间图形.
}

% ======================================================
% 第 69 天
% ======================================================
\setday{69}{重积分的换元法}

\oneproblem{
    求二重积分 $\displaystyle \iint_D (x-y)\,\dd x \dd y$，其中 $D=\{(x,y)|\ (x-1)^2+(y-1)^2\le2, y\ge x\}$.
}

\oneproblem{
    设有界区域 $D$ 是 $x^2+y^2=1$ 和直线 $y=x$ 及 $x$ 轴在第一象限围成的部分，计算二重积分 $I=\displaystyle \iint_D e^{(x+y)^2}(x^2-y^2)\,\dd x \dd y$.
}

\oneproblem{
    计算 $I=\displaystyle \iint_D \dfrac{3x}{y^2+xy^3}\,\dd x \dd y$，其中 $D$ 是由平面曲线 $xy=1, xy=3, y^2=x, y^2=3x$ 所围成的有界闭区域.
}

\oneproblem{
    计算 $I=\displaystyle \iint_D \dfrac{x^2-y^2}{\sqrt{x+y+3}}\,\dd x \dd y$，其中 $D: |x|+|y|\le1$.
}

\oneproblem{
    求由直线 $x+y=c, x+y=d, y=ax, y=bx\ (0<c<d, 0<a<b)$ 所围成的闭区域 $D$ 的面积.
}

\oneproblem{
    计算 $I=\displaystyle \iint_{x^2+y^2\le1} \left|\frac{x+y}{\sqrt{2}}-x^2-y^2\right|\,\dd x \dd y$.
}

\oneproblem{
    设 $\Omega: (x-2)^2+(y-3)^2+(z-4)^2\le1$，计算 $I=\displaystyle \iiint_{\Omega} (x^2+2y^2+3z^2)\,\dd x \dd ydz$.
}

\oneproblem{
    求由六个平面 $\begin{cases} a_1x+b_1y+c_1z=\pm h_1 \\ a_2x+b_2y+c_2z=\pm h_2 \\ a_3x+b_3y+c_3z=\pm h_3 \end{cases}$ 所围成的平行六面体的体积，其中 $A=\begin{vmatrix} a_1 & b_1 & c_1 \\ a_2 & b_2 & c_2 \\ a_3 & b_3 & c_3 \end{vmatrix}\neq0$ 且 $a_i,b_i,c_i,h_i\ (i=1,2,3)$ 为常数.
}

\oneproblem{
    设 $D=\{(x,y)|\ x^2+y^2\le1\}$，且 $a^2+b^2\neq0$，证明：
    \[
    \iint_D f(ax+by+c)\,\dd x \dd y = 2\int_{-1}^1 \sqrt{1-u^2} f\left(u\sqrt{a^2+b^2}+c\right)\,du.
    \]
}

\oneproblem{
    计算 $\displaystyle \iint_D \frac{(x+y)\ln\left(1+\frac{y}{x}\right)}{\sqrt{1-x-y}}\,\dd x \dd y$，其中区域 $D$ 由直线 $x+y=1$ 与两坐标轴所围三角形区域.
}

\oneproblem{
    某物体所在的空间区域为 $\Omega: x^2+y^2+2z^2\le x+y+2z$，密度函数为 $x^2+y^2+z^2$，求质量 $M=\displaystyle \iiint_{\Omega} (x^2+y^2+z^2)\,\dd x \dd ydz$.
}

\oneproblem{
    设 $f(x,y)$ 为 $\mathbb{R}^2$ 上的非负的连续函数，若 $I=\displaystyle \lim_{t\to+\infty} \iint_{x^2+y^2\le t^2} f(x,y)\,d\sigma$ 存在极限\\则称广义积分 $\displaystyle \iint_{\mathbb{R}^2} f(x,y)\,d\sigma$ 收敛于 $I$.
    \begin{enumerate}[label=(\arabic*)]
        \item 设 $f(x,y)$ 为 $\mathbb{R}^2$ 上的非负的连续函数，若 $\displaystyle \iint_{\mathbb{R}^2} f(x,y)\,d\sigma$ 收敛于 $I$\\证明极限 $\displaystyle \lim_{t\to+\infty} \iint_{-t\le x,y\le t} f(x,y)\,d\sigma$ 存在且收敛于 $I$.
        \item 设 $\displaystyle \iint_{\mathbb{R}^2} e^{ax^2+2bxy+cy^2}\,d\sigma$ 收敛于 $I$，实二次型 $ax^2+2bxy+cy^2$ 在正交变换下的标准二次型为 $\lambda_1u^2+\lambda_2v^2$。\\证明 $\lambda_1,\lambda_2$ 都小于零.
    \end{enumerate}
}

% ======================================================
% 第 70 天
% ======================================================
\setday{70}{一些具有特定结构的重积分的计算}

\oneproblem{
    已知函数 $f(x)=\begin{cases} e^x, & 0\le x\le1 \\ 0, & \text{其他} \end{cases}$，计算二重积分 $I=\displaystyle \int_{-\infty}^{+\infty} dx\int_{-\infty}^{+\infty} f(x)f(y-x)\,dy$.
}

\oneproblem{
    计算 $\displaystyle \int_{-\infty}^{+\infty}\int_{-\infty}^{+\infty} \min\{x,y\}e^{-(x^2+y^2)}\,\dd x \dd y$.
}

\oneproblem{
    已知 $\displaystyle \int_0^{+\infty} \frac{\sin x}{x}\,dx=\frac{\pi}{2}$，计算 $I=\displaystyle \int_0^{+\infty}\int_0^{+\infty} \frac{\sin x\sin(x+y)}{x(x+y)}\,\dd x \dd y$.
}

\oneproblem{
    计算二重积分 $\displaystyle \iint_D |x^2+y^2-1|\,d\sigma$，其中 $D=\{(x,y)|\ 0\le x\le1, 0\le y\le1\}$.
}

\oneproblem{
    求二重积分 $I=\displaystyle \iint_{x^2+y^2\le1} |x^2+y^2-x-y|\,\dd x \dd y$.
}

\oneproblem{
    求 $\displaystyle \iint_D \text{sgn}(xy-1)\,\dd x \dd y$，其中 $D=\{(x,y)|\ 0\le x\le2, 0\le y\le2\}$.
}

\oneproblem{
    计算二重积分 $\displaystyle \iint_D e^{\max\{x^2,y^2\}}\,\dd x \dd y$，其中 $D=\{(x,y)|\ 0\le x\le1, 0\le y\le1\}$.
}

\oneproblem{
    计算二重积分 $\displaystyle \iint_D xy[1+x^2+y^2]\,\dd x \dd y$，其中 $D=\{(x,y)|\ x^2+y^2\le\sqrt{2}, x\ge0, y\ge0\}$，$[1+x^2+y^2]$ 表示不超过 $1+x^2+y^2$ 的最大整数.
}

\oneproblem{
    计算二重积分 $\displaystyle \iint_D f(x,y)\,d\sigma$，其中二元函数 $f(x,y)=\begin{cases} x^2, & |x|+|y|\le1 \\ \dfrac{1}{\sqrt{x^2+y^2}}, & 1<|x|+|y|\le2 \end{cases}$，积分区域 $D=\{(x,y)|\ |x|+|y|\le2\}$.
}

\oneproblem{
    试用投影法、截面法、球面坐标方法和一种换元法计算三重积分 $I=\displaystyle \iiint_{\Omega} \dfrac{dV}{\rho^2}$，其中 $\rho$ 是点 $(x,y,z)$ 到 $x$ 轴的距离，$\Omega$ 为一棱台，其六个顶点分别为 $A(0,0,1),B(0,1,1),C(1,1,1),D(0,0,2),E(0,2,2),F(2,2,2)$.
}